% =====================================================================
%  PRO3600 --- Rapport de projet : Plateforme web de backtesting
%  Compilation : pdflatex rapport.tex  (x2 pour la table des matières)
%  ou : latexmk -pdf rapport.tex
%  Compatible Overleaf (moteur pdfLaTeX).
% =====================================================================
\documentclass[11pt,a4paper]{article}

\usepackage[utf8]{inputenc}
\usepackage[T1]{fontenc}
\usepackage[french]{babel}
\usepackage{lmodern}
\usepackage{amsmath}
\usepackage{textcomp}
\usepackage[a4paper,margin=2.4cm]{geometry}
\usepackage{graphicx}
\usepackage{float}
\usepackage{xcolor}
\usepackage{booktabs}
\usepackage{longtable}
\usepackage{array}
\usepackage{enumitem}
\usepackage{listings}
\usepackage{fancyhdr}
\usepackage{titlesec}
\usepackage{tikz}
\usetikzlibrary{positioning,arrows.meta,fit,calc}
\usepackage[hidelinks]{hyperref}

% ---- Couleurs ----
\definecolor{codebg}{RGB}{245,245,245}
\definecolor{codekw}{RGB}{0,90,160}
\definecolor{codecmt}{RGB}{100,120,100}
\definecolor{codestr}{RGB}{160,60,40}
\definecolor{accent}{RGB}{30,90,150}

% ---- Listings (code) ----
\lstset{
  basicstyle=\ttfamily\footnotesize,
  backgroundcolor=\color{codebg},
  keywordstyle=\color{codekw}\bfseries,
  commentstyle=\color{codecmt}\itshape,
  stringstyle=\color{codestr},
  numbers=left, numberstyle=\tiny\color{gray}, numbersep=8pt,
  frame=single, rulecolor=\color{gray!40},
  breaklines=true, showstringspaces=false,
  tabsize=2, captionpos=b, xleftmargin=14pt,
}
\lstdefinestyle{py}{language=Python}
\lstdefinestyle{ts}{language=Java} % proche TypeScript pour la coloration
\lstdefinestyle{sh}{language=bash}

% ---- En-têtes ----
\pagestyle{fancy}
\fancyhf{}
\fancyhead[L]{\small PRO3600 --- Backtesting}
\fancyhead[R]{\small\leftmark}
\fancyfoot[C]{\thepage}
\renewcommand{\headrulewidth}{0.4pt}

\titleformat{\section}{\Large\bfseries\color{accent}}{\thesection}{1em}{}
\titleformat{\subsection}{\large\bfseries}{\thesubsection}{1em}{}

\newcommand{\code}[1]{\texttt{#1}}

% ---- Images : placeholders intelligents ----
% Déposez les images dans le sous-dossier images/. Tant qu'un fichier est absent,
% une case grise légendée s'affiche (le document compile dans tous les cas) ;
% dès que le fichier existe, il est inséré automatiquement.
\graphicspath{{images/}}
% \figph{fichier}{largeur}{hauteur-case}{légende}
\newcommand{\figph}[4]{%
  \begin{figure}[H]\centering
  \IfFileExists{images/#1}{\includegraphics[width=#2]{#1}}{%
    {\setlength{\fboxsep}{0pt}%
     \fcolorbox{gray!55}{gray!10}{\parbox[c][#3][c]{#2}{\centering
       \footnotesize\itshape\color{gray!80}[\,image à insérer\,]\\[3pt]\texttt{images/#1}}}}}%
  \caption{#4}
  \end{figure}}
% Logo (sans flottant, pour la page de garde)
\newcommand{\logoph}{%
  \IfFileExists{images/logo-tsp.png}{\includegraphics[width=3.8cm]{logo-tsp.png}}{%
    {\setlength{\fboxsep}{0pt}%
     \fcolorbox{gray!55}{gray!10}{\parbox[c][2cm][c]{3.8cm}{\centering
       \footnotesize\itshape\color{gray!80}[\,logo TSP\,]\\[2pt]\texttt{images/logo-tsp.png}}}}}}

% Évite les débordements de ligne (overfull hbox) dans les paragraphes
\setlength{\emergencystretch}{3em}
\sloppy

% =====================================================================
\begin{document}

% --------------------- PAGE DE GARDE ---------------------
\begin{titlepage}
\centering
\vspace*{1.0cm}
\logoph\par
\vspace{1.0cm}
{\scshape\Large Télécom SudParis --- Cycle Ingénieur\par}
\vspace{0.4cm}
{\scshape\large Module PRO 3600 --- Projet de Développement Logiciel\par}
\vspace{2.5cm}
{\Huge\bfseries RAPPORT DE PROJET\par}
\vspace{0.5cm}
{\LARGE\bfseries Plateforme web de backtesting algorithmique\par}
\vspace{2.5cm}
{\large\textbf{Équipe projet :}\par}
\vspace{0.3cm}
{\large Adrien \textsc{Etienne} \quad Neil \textsc{Kiuchi} \quad Elio \textsc{Bonnenfant}\par}
\vspace{0.2cm}
{\large Robinson \textsc{Cerceau} \quad Mattéo \textsc{Machnik-Hauchart}\par}
\vspace{1.5cm}
{\large\textbf{Encadrant du projet :}\par}
\vspace{0.3cm}
{\large M. \textsc{Romero}\par}
\vfill
{\large\textbf{Version du document :} 3.0 (Livrable final : logiciel, tests et rapport)\par}
{\large\textbf{Date :} \today\par}
\end{titlepage}

% --------------------- TABLE DES MATIÈRES ---------------------
\tableofcontents
\newpage

% =====================================================================
\section{Introduction}
% =====================================================================
Ce document constitue le rapport final du projet PRO3600 mené par notre équipe.
Il prolonge et enrichit le pré-rapport (version~1.0) en y intégrant l'histoire
complète du projet, la conception détaillée du logiciel effectivement livré, les
résultats de la campagne de tests, un bilan objectif/réalisation et un manuel
utilisateur. Conformément aux attendus du livrable~3, il vise à permettre à une
personne ignorant tout du projet d'en comprendre les tenants et aboutissants, à
un développeur de maintenir le logiciel, et à un utilisateur de se servir de
l'application.

\medskip
Le logiciel livré est une \textbf{plateforme web de backtesting de stratégies de
trading} sur données financières réelles. Elle est composée d'une API
\textbf{FastAPI} (Python) et d'une application monopage \textbf{Angular~21}, et
permet à un utilisateur de tester des stratégies d'analyse technique (MACD, RSI,
croisement de moyennes mobiles) sur des indices boursiers, d'en visualiser les
performances et de conserver un historique de ses analyses.

% =====================================================================
\section{Cadrage du projet et périmètre}
% =====================================================================
\subsection{Motivation et objectifs}
En finance quantitative, il est indispensable de tester rigoureusement une
stratégie d'investissement sur des données historiques avant de l'appliquer sur
les marchés réels avec de véritables capitaux. Ce processus de simulation,
appelé \emph{backtesting}, permet de valider des intuitions mathématiques et
d'évaluer la robustesse d'un modèle face aux conditions passées. L'objectif du
projet est de concevoir une plateforme web interactive dédiée à cette pratique,
agissant comme un laboratoire virtuel d'analyse de modèles d'investissement.

\medskip
Les solutions de backtesting existantes sont souvent complexes ou nécessitent des
compétences avancées en programmation (scripts, API). La valeur ajoutée de notre
projet est d'offrir une alternative \textbf{ergonomique et visuelle}, rendant
l'analyse quantitative accessible directement depuis un navigateur web :
l'application évalue automatiquement performances et risques, et restitue les
résultats sous forme de graphiques et de métriques financières standards.

\subsection{Définition du Produit Minimum Viable (MVP)}
L'objectif du premier livrable était de valider l'architecture globale
(communication Front-End / Back-End) et de démontrer la faisabilité des calculs.
Le MVP visait à permettre à un utilisateur de charger un jeu de données, d'y
appliquer une stratégie de base et de visualiser le résultat financier. Le
périmètre du MVP comprenait :

\begin{enumerate}[label=\arabic*)]
  \item \textbf{Gestion des données} : chargement d'un historique financier local
        (CSV, p.~ex. le S\&P~500 sur 5 ans).
  \item \textbf{Moteur de backtesting} : implémentation de plusieurs stratégies
        d'analyse technique standard (MA, MA Crossover, MACD, RSI) et simulation
        de l'exécution des ordres d'achat/vente.
  \item \textbf{Évaluation et métriques} : calcul du rendement cumulatif de la
        stratégie sur la période.
  \item \textbf{Interface utilisateur} : une interface web, un bouton de
        déclenchement du backtest, un graphique du prix de l'actif et l'affichage
        de la performance finale.
\end{enumerate}

Étaient explicitement \emph{exclus} du MVP : la connexion à des API publiques
temps réel, une base de données pour les historiques utilisateurs, des stratégies
complexes multiples et les métriques de risque avancées (Sharpe, volatilité).

\subsection{Évolution du périmètre : au-delà du MVP}
\label{sec:evolution}
Une fois le MVP validé, l'équipe a poursuivi le développement bien au-delà du
périmètre initial. Plusieurs fonctionnalités classées « après-MVP » dans le
pré-rapport ont en réalité été \textbf{réalisées} :
\begin{itemize}[leftmargin=1.4em]
  \item connexion à une \textbf{API publique temps réel} (Yahoo Finance via
        \code{yfinance}) à la place du CSV statique ;
  \item \textbf{base de données} (Supabase PostgreSQL) pour persister l'historique
        des backtests par utilisateur ;
  \item \textbf{métriques de risque avancées} : ratio de Sharpe annualisé,
        \emph{maximum drawdown}, taux de réussite, nombre de transactions ;
  \item \textbf{authentification} via Google OAuth et un mode démo sans compte.
\end{itemize}
Ce dépassement du périmètre est analysé dans le bilan (section~\ref{sec:bilan}).

\subsection{Stack technologique}
\begin{center}
\begin{tabular}{@{}ll@{}}
\toprule
\textbf{Couche} & \textbf{Technologies} \\
\midrule
Back-End  & Python 3.11, FastAPI, Pandas, NumPy, \code{yfinance}, SQLAlchemy~2.0 \\
Base de données & SQLite (développement) / PostgreSQL Supabase (production) \\
Authentification & Authlib (Google OAuth/OIDC), sessions signées Starlette \\
Front-End & Angular~21 (standalone), Angular Material, \code{lightweight-charts} \\
Tests     & pytest + coverage (back) ; Vitest + Angular TestBed (front) \\
Outils    & Git \& GitHub, Discord, Google Drive \\
\bottomrule
\end{tabular}
\end{center}

% =====================================================================
\section{Organisation de l'équipe}
% =====================================================================
Afin de paralléliser le travail et de garantir l'efficacité du développement,
l'équipe a défini des pôles de compétences (non figés).

\subsection{Attribution des rôles et pôles}
\begin{center}
\begin{tabular}{@{}p{4.2cm}p{4.4cm}p{5.6cm}@{}}
\toprule
\textbf{Pôle} & \textbf{Responsables} & \textbf{Périmètre} \\
\midrule
Frontend & Robinson Cerceau, Elio Bonnenfant & SPA Angular, interface de
paramétrage, intégration des graphiques, appels API. \\
Backend (architecture \& API) & Neil Kiuchi, Mattéo Machnik-Hauchart &
Architecture serveur, endpoints, gestion des requêtes, sécurisation et erreurs. \\
Data & Elio Bonnenfant & Extraction/nettoyage des historiques de prix, format
OHLCV, intégration de l'API publique. \\
Backtesting \& stratégies & Adrien Etienne & Simulateur de trading, algorithmes
d'analyse technique, génération des signaux, métriques d'évaluation. \\
\bottomrule
\end{tabular}
\end{center}

\subsection{Méthodologie de travail et outils collaboratifs}
Nous avons adopté une approche \textbf{Agile} : le développement est découpé en
cycles itératifs (sprints), permettant de rester flexibles, d'adapter les
objectifs aux difficultés rencontrées et de livrer de manière incrémentale. Un
\textbf{diagramme de Gantt} complète cette approche pour une vision macroscopique
et le repérage des jalons (livrables).

\medskip
L'environnement collaboratif s'articule autour de trois outils :
\begin{itemize}[leftmargin=1.4em]
  \item \textbf{Discord} --- communication temps réel (salons thématiques, points
        d'avancement, réunions vocales) ;
  \item \textbf{GitHub} --- gestion de version du code, suivi des modifications,
        travail simultané par branches (\code{features}, \code{fix}) ;
  \item \textbf{Google Drive} --- centralisation documentaire (cahiers des charges,
        comptes rendus, recherches bibliographiques, rédaction des rapports).
\end{itemize}

% =====================================================================
\section{Gestion et histoire du projet}
% =====================================================================
\subsection{Planification globale et jalons}
Un plan de charge sous forme de tableur a recensé l'ensemble des tâches, estimé
les temps de réalisation et réparti les responsabilités ; mis à jour
régulièrement, il a permis de comparer estimations et temps réellement investi.
Le diagramme de Gantt prévisionnel a servi de feuille de route. Les jalons
principaux sont rappelés ci-dessous.

\begin{center}
\begin{tabular}{@{}lll@{}}
\toprule
\textbf{Jalon} & \textbf{Échéance} & \textbf{Contenu} \\
\midrule
Livrable 1 & 23 février & Définition du MVP, planning du sprint~1 \\
Livrable 2 & 27 mars    & Itérations intermédiaires \\
Livrable 3 & 15 mai     & Logiciel, tests et rapport final \\
Soutenance & 11--12 juin & Présentation et démonstration \\
\bottomrule
\end{tabular}
\end{center}

\medskip
\noindent Phases planifiées (Gantt) : choix de l'architecture (janv.--23 fév.),
choix back/front (15 janv.--23 fév.), MVP (22 fév.--8 mars), développement frontend
(22 fév.--15 mai), développement backend avec données temps réel (22 fév.--24 avr.),
modélisation des tables de la base (3 mars--7 mai), développement du moteur de
backtesting et stratégies (10 mars--7 mai), phase de test (10 avr.--15 mai),
préparation de la soutenance (15 mai--11 juin).

\figph{gantt.png}{0.95\linewidth}{6.5cm}{Diagramme de Gantt prévisionnel du projet (jalons et phases de développement).}

\subsection{Stratégie de découpage en sprints}
Conformément à l'approche Agile, le développement a été découpé en itérations
courtes dotées d'objectifs précis. Le \textbf{sprint~1} (MVP, échéance 9~mars)
visait à mettre en place l'environnement de travail et à valider la communication
de bout en bout entre Angular et FastAPI sur un jeu de données.

\paragraph{Carnet du sprint 1 (extrait).}
\begin{itemize}[leftmargin=1.4em]
  \item \textbf{Pôle Architecture} : création du dépôt GitHub, arborescence
        \code{/backend} et \code{/frontend}, \code{.gitignore}, README de lancement.
  \item \textbf{Pôle Back-End \& Data} : initialisation de FastAPI et d'une route
        de test, lecture d'un CSV avec Pandas, route renvoyant les données en JSON.
  \item \textbf{Pôle Front-End} : initialisation du projet Angular, structure des
        pages, intégration d'une librairie de graphiques, appel \code{fetch} vers
        le backend.
\end{itemize}
Les sprints suivants ont porté sur le remplacement du CSV par l'API temps réel,
la modélisation de la base de données, le moteur de backtesting et les stratégies,
l'authentification, puis la phase de tests et de stabilisation.

\subsection{Choix technologiques et justifications}
\begin{itemize}[leftmargin=1.4em]
  \item \textbf{FastAPI} : framework Python moderne, performant et auto-documenté
        (OpenAPI/Swagger), idéal pour exposer rapidement une API typée avec
        validation des données (Pydantic).
  \item \textbf{Pandas / NumPy} : manipulation vectorisée des séries temporelles
        financières, calculs d'indicateurs et de rendements concis et efficaces.
  \item \textbf{yfinance} : accès gratuit aux données de marché Yahoo Finance,
        retenu pour passer du CSV statique à des données réelles sans coût d'API.
  \item \textbf{Angular 21 (standalone)} : framework structurant, typé
        (TypeScript), avec un système de \emph{signals} pour la réactivité et
        Angular Material pour une UI cohérente.
  \item \textbf{lightweight-charts} (TradingView) : librairie de graphiques
        financiers performante (chandeliers, courbes) adaptée aux séries longues.
  \item \textbf{SQLAlchemy 2.0 + PostgreSQL (Supabase)} : ORM moderne et base
        managée pour persister les historiques ; SQLite en développement local.
\end{itemize}

\subsection{Problèmes rencontrés et solutions}
\begin{itemize}[leftmargin=1.4em]
  \item \textbf{Colonnes MultiIndex de yfinance} : pour un seul ticker,
        \code{yfinance} renvoie des colonnes hiérarchiques. Solution : aplatissement
        via \code{df.columns.get\_level\_values(0)} et renommage de l'index en
        \code{Time}.
  \item \textbf{Biais de look-ahead} : un backtest naïf utiliserait le signal du
        jour~$t$ pour le rendement du jour~$t$, ce qui « triche ». Solution :
        décalage du signal d'une barre (\code{signal.shift(1)}) dans le moteur, de
        sorte que la position décidée en $t$ s'applique au rendement de $t{+}1$.
  \item \textbf{CORS Front/Back} : l'application Angular (port 4200 / Vercel)
        appelle l'API (port 8000 / Render) depuis une autre origine. Solution :
        configuration du \code{CORSMiddleware} avec la liste blanche des origines.
  \item \textbf{Authentification Google} : mise en place du flux OAuth/OIDC
        (Authlib) avec redirection vers le frontend après connexion.
  \item \textbf{Angular 21 « zoneless »} : la nouvelle version n'utilise plus
        \code{zone.js}, ce qui a nécessité d'adapter la réactivité (signals,
        \code{effect}) et la configuration des tests.
  \item \textbf{Outillage de test} : le \emph{scaffolding} de tests frontend
        n'était pas câblé. Solution : ajout d'un \code{tsconfig.spec.json} incluant
        les \code{*.spec.ts} et des types \code{vitest/globals} (voir
        section~\ref{sec:tests}).
\end{itemize}

% =====================================================================
\section{Conception détaillée}
% =====================================================================
\subsection{Architecture générale}
La plateforme repose sur une architecture client/serveur découplée, organisée en
trois couches aux responsabilités distinctes :

\begin{description}[leftmargin=1.4em,style=nextline]
  \item[Couche présentation (client).] Une application monopage (SPA) Angular~21,
        déployée sur Vercel. Elle gère l'affichage et les interactions : les pages
        (connexion, tableau de bord, données, backtests, mode avancé, résultats,
        paramètres), les services d'accès aux données et à l'authentification, et
        les composants de graphiques financiers.
  \item[Couche application (serveur).] Une API REST FastAPI, déployée sur Render.
        Elle reçoit les requêtes du client, récupère les cours de marché, exécute
        les stratégies et le moteur de backtest, calcule les métriques et gère
        l'authentification.
  \item[Couche persistance (données).] Une base de données relationnelle accédée
        via l'ORM SQLAlchemy~2.0 : SQLite en développement local, PostgreSQL
        (hébergé sur Supabase) en production.
\end{description}

\noindent Les couches communiquent par des contrats clairs. Le client appelle
l'API en \textbf{HTTP/JSON} (style REST). Le serveur s'appuie sur deux services
externes : \textbf{Yahoo Finance} (via la librairie \code{yfinance}) pour les
cours, avec une \textbf{mise en cache} de 24~heures en base afin de limiter les
appels réseau, et \textbf{Google} (OAuth\,2.0 / OIDC) pour l'authentification. La
persistance des utilisateurs, des backtests et du cache passe par SQLAlchemy. La
figure~\ref{fig:archi} en donne une vue d'ensemble, et la figure~\ref{fig:bdd}
détaille le schéma de la base de données.

\begin{figure}[H]
\centering
\resizebox{\textwidth}{!}{%
\begin{tikzpicture}[
  tier/.style={rectangle, rounded corners, draw=accent, line width=1pt,
               fill=accent!5, align=left, inner sep=10pt, text width=128mm},
  ext/.style={rectangle, rounded corners, draw=gray!70, fill=gray!10,
              align=center, inner sep=8pt, text width=34mm, font=\small},
  lbl/.style={font=\small\itshape, fill=white, inner sep=2pt},
  flow/.style={-{Latex[length=3mm]}, line width=1.1pt, draw=gray!75},
  bidir/.style={{Latex[length=3mm]}-{Latex[length=3mm]}, line width=1.1pt, draw=gray!75},
]
\node[tier](client){%
  \textbf{Couche présentation --- Client (navigateur web)}\hfill\textit{déploiement : Vercel}\\[3pt]
  Application monopage \textbf{Angular 21} (standalone, zoneless)\\
  \textbf{Pages} : login, dashboard, data, backtests, advanced, results, settings\\
  \textbf{Services} : AuthService, DataService, BacktestHistoryService\\
  \textbf{Graphiques} : market-chart, equity-chart \textit{(lightweight-charts)}};

\node[tier, below=24mm of client](api){%
  \textbf{Couche application --- API REST}\hfill\textit{déploiement : Render}\\[3pt]
  Serveur \textbf{FastAPI} (\texttt{main:app}, lancé par Uvicorn)\\
  \textbf{Endpoints} : /api/data, /api/backtest, /api/backtest/upload, /api/history,\\
  \hspace*{1.4em}/api/strategies, /api/auth/*\\
  \textbf{Modules} : moteur \texttt{backtest.py}, \texttt{strategies/} (MACD, RSI, MA),\\
  \hspace*{1.4em}\texttt{metrics.py} (Sharpe, Sortino, profit factor\dots), \texttt{validator.py},\\
  \hspace*{1.4em}\texttt{cache.py} (cache 24\,h), \texttt{auth.py} (OAuth)};

\node[tier, below=24mm of api](data){%
  \textbf{Couche persistance --- Données}\\[3pt]
  ORM \textbf{SQLAlchemy 2.0} $\rightarrow$ \textbf{SQLite} (dév.) / \textbf{PostgreSQL Supabase} (prod.)\\
  \textbf{Tables} : users, backtests, saved\_strategies, watchlist, market\_data\_cache};

\node[ext, left=14mm of api](yahoo){\textbf{Yahoo Finance}\\\textit{librairie yfinance}\\(cours OHLCV)};
\node[ext, right=14mm of api](google){\textbf{Google}\\\textit{OAuth 2.0 / OIDC}\\(authentification)};

\draw[flow] (client) -- node[lbl]{HTTP / JSON (REST)} (api);
\draw[flow] (api) -- node[lbl]{SQLAlchemy ORM} (data);
\draw[bidir] (api) -- node[lbl]{OHLCV} (yahoo);
\draw[bidir] (api) -- node[lbl]{OAuth} (google);
\end{tikzpicture}}
\caption{Architecture globale de la plateforme : trois couches découplées
(présentation Angular, application FastAPI, persistance SQLAlchemy) et deux
services externes (Yahoo Finance pour les cours, Google pour l'authentification).}
\label{fig:archi}
\end{figure}

\subsection{Backend : structure et API}
Le backend est organisé en modules à responsabilité unique :
\begin{center}
\begin{tabular}{@{}l p{10.8cm}@{}}
\toprule
\textbf{Fichier} & \textbf{Rôle} \\
\midrule
\code{main.py}        & Point d'entrée, middlewares, endpoints backtest/historique/stratégies/upload \\
\code{data.py}        & Routeur \code{/api/data} : récupération OHLCV pour les graphiques \\
\code{backtest.py}    & Moteur de backtest vectorisé (rendements, métriques) \\
\code{strategies/}    & Stratégies (MACD, RSI, MA crossover) + registre \\
\code{metrics.py}     & Métriques du \emph{mode avancé} (Sortino, profit factor, espérance\dots) sur trades importés \\
\code{validator.py}   & Validation des backtests importés (mode avancé) \\
\code{cache.py}       & Cache des cours de marché en base (TTL 24\,h) \\
\code{models.py}      & Modèles ORM (User, Backtest, SavedStrategy, Watchlist, MarketDataCache) \\
\code{database.py}    & Configuration SQLAlchemy (engine, session, \code{get\_db}) \\
\code{auth.py}        & Routeur \code{/api/auth} : Google OAuth, gestion de compte \\
\bottomrule
\end{tabular}
\end{center}

\paragraph{Endpoints de l'API.}
\begin{center}
\begin{tabular}{@{}lll@{}}
\toprule
\textbf{Méthode} & \textbf{Route} & \textbf{Description} \\
\midrule
GET    & \code{/api/data/\{ticker\}?period=} & Données OHLCV (yfinance) \\
POST   & \code{/api/backtest}               & Lance un backtest \\
POST   & \code{/api/backtest/upload}        & Backtest en mode avancé (import de trades JSON) \\
GET    & \code{/api/history/\{user\_id\}}     & Historique d'un utilisateur \\
DELETE & \code{/api/history/\{id\}}           & Supprime un backtest \\
GET    & \code{/api/strategies}             & Liste des stratégies \\
GET    & \code{/api/auth/login}             & Redirection vers Google \\
GET    & \code{/api/auth/callback}          & Retour OAuth, création/MAJ utilisateur \\
PATCH  & \code{/api/auth/users/\{id\}}        & Renommage d'utilisateur \\
DELETE & \code{/api/auth/users/\{id\}}        & Suppression de compte \\
\bottomrule
\end{tabular}
\end{center}

\paragraph{Modèle de données.}
La base comporte cinq tables (figure~\ref{fig:bdd}). La table \code{users} et la
table \code{backtests} constituent le c\oe{}ur persistant : un enregistrement de
backtest référence son utilisateur (clé étrangère \code{user\_id}) et stocke le
contexte (ticker, stratégie, période, paramètres, capital, stop-loss) ainsi que
les métriques calculées (rendement stratégie/marché, Sharpe, drawdown, nombre de
transactions, taux de réussite). La table \code{market\_data\_cache} met en cache
les cours récupérés auprès de Yahoo Finance, avec une contrainte d'unicité sur le
couple \code{(ticker, period)} et une durée de vie de 24~heures. Deux tables
(\code{saved\_strategies}, \code{watchlist}) sont préparées pour des évolutions.

\begin{figure}[H]
\centering
\resizebox{0.96\textwidth}{!}{%
\begin{tikzpicture}[
  entity/.style={rectangle, draw=accent, line width=0.8pt, inner sep=0pt, font=\small},
  rel/.style={line width=1pt, draw=accent!70},
  card/.style={font=\scriptsize\bfseries, fill=white, inner sep=1.5pt},
  verb/.style={font=\small\itshape, fill=white, inner sep=2pt},
]
% --- Entité parente ---
\node[entity](users){\begin{tabular}{@{\,}l@{\,}}
  \textbf{users}\\\hline
  \underline{id} : int (PK)\\ username : str (UNIQUE, NN)\\ email : str (UNIQUE, NN)\\
  password\_hash : str (NN)\\ created\_at : datetime\\
\end{tabular}};

% --- Entités filles : référencent users via user\_id ---
\node[entity, right=54mm of users, yshift=40mm](bt){\begin{tabular}{@{\,}l@{\,}}
  \textbf{backtests}\\\hline
  \underline{id} : int (PK)\\ \textit{user\_id} : int (FK)\\
  ticker, strategy, period : str\\ params : JSON\\
  capital\_initial, stop\_loss : float\\ total\_return\_strat / \_market : float\\
  sharpe\_ratio, max\_drawdown : float\\ n\_trades : int, win\_rate : float\\
  created\_at : datetime\\
\end{tabular}};

\node[entity, right=54mm of users](ss){\begin{tabular}{@{\,}l@{\,}}
  \textbf{saved\_strategies}\\\hline
  \underline{id} : int (PK)\\ \textit{user\_id} : int (FK)\\
  name, strategy : str\\ params : JSON\\ created\_at : datetime\\
\end{tabular}};

\node[entity, right=54mm of users, yshift=-46mm](wl){\begin{tabular}{@{\,}l@{\,}}
  \textbf{watchlist}\\\hline
  \underline{id} : int (PK)\\ \textit{user\_id} : int (FK)\\
  ticker : str\\ added\_at : datetime\\
\end{tabular}};

% --- Entité autonome (aucune relation) ---
\node[entity, below=18mm of users](cache){\begin{tabular}{@{\,}l@{\,}}
  \textbf{market\_data\_cache}\\\hline
  \underline{id} : int (PK)\\ ticker : str, period : str\\
  \textit{(UNIQUE ticker, period)}\\ data : JSON\\ fetched\_at : datetime\\
  \textit{(table autonome, aucune FK)}\\
\end{tabular}};

% --- Relations 1:N : un user possède 0..N lignes ; chaque ligne -> 1 user ---
\draw[rel] (users.east) to[out=35,in=180] (bt.west)
  node[card, pos=0.12]{1} node[verb, pos=0.56]{exécute} node[card, pos=0.93]{0..N};
\draw[rel] (users.east) to[out=0,in=180] (ss.west)
  node[card, pos=0.12]{1} node[verb, pos=0.50]{enregistre} node[card, pos=0.90]{0..N};
\draw[rel] (users.east) to[out=-35,in=180] (wl.west)
  node[card, pos=0.12]{1} node[verb, pos=0.56]{suit} node[card, pos=0.93]{0..N};
\end{tikzpicture}}
\caption{Schéma relationnel de la base de données (cardinalités 1:N). Un
\texttt{users} possède 0..N lignes dans \texttt{backtests}, \texttt{saved\_strategies}
et \texttt{watchlist} via la clé étrangère \texttt{user\_id} ; réciproquement, chaque
ligne appartient à exactement un utilisateur. \texttt{market\_data\_cache} est
autonome (cache des cours, contrainte d'unicité sur \texttt{(ticker, period)}).
\textit{PK : clé primaire \textbullet{} FK : clé étrangère \textbullet{} NN : non nul.}}
\label{fig:bdd}
\end{figure}

\paragraph{Moteur de backtest.}
Le moteur (\code{backtest.py}) est entièrement vectorisé. À partir d'un
DataFrame contenant la colonne de signaux, il calcule :
\begin{lstlisting}[style=py]
market_return = Close.pct_change()
strat_return  = market_return * signal.shift(1)   # decalage = anti look-ahead
cumulative    = (1 + return).cumprod() - 1
\end{lstlisting}
Les métriques produites par le moteur standard sont : rendement cumulé stratégie
et marché, ratio de Sharpe annualisé ($\sqrt{252}$), \emph{maximum drawdown},
nombre de transactions (changements de signal) et taux de réussite (proportion de
barres en position gagnantes). Le \emph{mode avancé} (import de trades, calculé
par le module \code{metrics.py}) produit en complément des métriques
supplémentaires : ratio de \textbf{Sortino}, \textbf{profit factor}, espérance par
trade, volatilité annualisée, etc.

\paragraph{Stratégies.}
Toutes les stratégies héritent d'une interface commune \code{BaseStrategy} et
implémentent \code{compute\_signals(df) -> df}, qui ajoute une colonne
\code{signal} ($+1$ long, $-1$ court, $0$ neutre) et des colonnes d'indicateurs.
Un registre (\code{STRATEGIES\_REGISTRY}) associe une clé textuelle à chaque
classe, consommée par la route \code{/api/backtest}.
\begin{center}
\begin{tabular}{@{}ll p{6.6cm}@{}}
\toprule
\textbf{Stratégie} & \textbf{Paramètres} & \textbf{Logique de signal} \\
\midrule
MACD          & fast=12, slow=26, signal=9 & $+1$ si MACD\_line $>$ signal, sinon $-1$ \\
RSI           & length=14, ob=70, os=30    & $+1$ si RSI $<$ os, $-1$ si RSI $>$ ob, sinon $0$ \\
MA crossover  & fast=10, slow=30           & $+1$ si MA\_fast $>$ MA\_slow, sinon $-1$ \\
Bollinger     & length=20, num\_std=2      & $+1$ sous la bande basse, $-1$ au-dessus de la bande haute, sinon $0$ \\
Momentum      & length=10                  & $+1$ si variation sur $N$ barres $>0$, $-1$ si $<0$ \\
Donchian      & length=20                  & $+1$/$-1$ à la cassure du plus haut/bas des $N$ barres précédentes \\
\bottomrule
\end{tabular}
\end{center}

\subsection{Frontend : structure et logique}
L'application Angular est \emph{standalone} (sans NgModule). Le composant racine
affiche une barre de navigation (\code{navbar}) qui héberge le
\code{<router-outlet>}. Les pages sont chargées paresseusement (lazy loading) et
protégées par un garde de route.

\paragraph{Routage.} \code{login} et \code{callback} sont publiques ; les pages
\code{dashboard}, \code{data}, \code{backtests}, \code{results}, \code{settings}
sont protégées par \code{authGuard}.

\paragraph{Services.}
\begin{itemize}[leftmargin=1.4em]
  \item \code{DataService} --- tous les appels REST (données, backtest, historique,
        compte) ; l'URL de base provient de \code{environment.apiUrl}.
  \item \code{AuthService} --- identité stockée en \code{localStorage}, mode démo,
        état d'authentification.
  \item \code{BacktestHistoryService} --- historique local (mode démo), tampon
        circulaire des 5 derniers résultats.
\end{itemize}

\paragraph{Composants graphiques.} \code{market-chart} affiche les cours
(chandeliers ou ligne selon la période) et \code{equity-chart} la courbe
d'équité (stratégie vs marché, rebasée en base~100), via \code{lightweight-charts}.

\subsection{Flux de bout en bout d'un backtest}
\begin{enumerate}[leftmargin=1.6em]
  \item L'utilisateur choisit ticker, période, stratégie et paramètres sur la page
        \code{Backtests} et clique sur « Lancer ».
  \item \code{DataService.runBacktest()} envoie un \code{POST /api/backtest}.
  \item L'API valide la stratégie, récupère les données, instancie la stratégie,
        calcule les signaux puis exécute le moteur de backtest.
  \item La réponse (\code{stats}, \code{equity\_curve}, \code{signals}) est
        renvoyée ; si l'utilisateur est connecté, le résultat est persisté.
  \item Le frontend ajoute le résultat à l'historique local et navigue vers la page
        \code{Results} qui affiche métriques, courbe d'équité et tableau de signaux.
\end{enumerate}

% =====================================================================
\section{Réalisation : fonctionnalités livrées}
% =====================================================================
\begin{itemize}[leftmargin=1.4em]
  \item Récupération de données de marché temps réel (S\&P~500, NASDAQ, NQ) sur
        plusieurs périodes (1J, 1M, 1A, 5A) et visualisation graphique.
  \item Lancement de backtests sur six stratégies paramétrables (MACD, RSI, MA
        crossover, Bollinger, Momentum, Donchian) avec calcul de métriques.
  \item Page de résultats : métriques, courbe d'équité interactive, tableau des
        signaux, navigation dans l'historique récent.
  \item Tableau de bord synthétique (nombre de backtests, meilleur rendement,
        runs récents).
  \item Authentification Google OAuth + mode démo ; persistance de l'historique en
        base pour les utilisateurs connectés.
  \item Page de paramètres (nom d'utilisateur, préférences d'interface, suppression
        de compte).
\end{itemize}

\medskip
\noindent Les captures ci-dessous illustrent les principales interfaces livrées.

\figph{screenshot-dashboard.png}{0.85\linewidth}{6cm}{Tableau de bord : synthèse des backtests et meilleur rendement.}

\figph{screenshot-backtests.png}{0.85\linewidth}{6cm}{Page Backtests : sélection de l'indice, de la période, de la stratégie et de ses paramètres.}

\figph{screenshot-results.png}{0.85\linewidth}{6cm}{Page Résultats : métriques, courbe d'équité (stratégie vs marché) et tableau des signaux.}

\figph{screenshot-data.png}{0.85\linewidth}{6cm}{Page Données de marché : graphique des cours (chandeliers / ligne) via lightweight-charts.}

% =====================================================================
\section{Tests et validation}
\label{sec:tests}
% =====================================================================
Nous avons écrit deux types de tests automatisés, classiques en développement
logiciel :
\begin{itemize}[leftmargin=1.4em,nosep]
  \item des \textbf{tests unitaires}, qui vérifient une fonction à la fois (par
        exemple : une stratégie produit bien un signal cohérent, ou un calcul de
        rendement donne le bon résultat) ;
  \item des \textbf{tests d'intégration}, qui vérifient que l'API répond
        correctement de bout en bout (lancer un backtest, lire ou supprimer un
        historique, et renvoyer une erreur quand la requête est invalide).
\end{itemize}
Ils sont écrits avec \textbf{pytest} côté backend (Python) et avec \textbf{Vitest}
côté frontend (Angular). Nous avons complété par quelques essais manuels, en
utilisant l'application comme un utilisateur (connexion, lancement d'un backtest,
lecture des résultats).

\subsection{Résultat}
\textbf{L'ensemble des tests passe} : 99 tests automatisés sont au vert.
\begin{center}
\begin{tabular}{@{}lcc@{}}
\toprule
\textbf{Périmètre} & \textbf{Nombre de tests} & \textbf{Résultat} \\
\midrule
Backend (pytest)  & 58 & tous au vert \\
Frontend (Vitest) & 41 & tous au vert \\
\textbf{Total}    & \textbf{99} & \textbf{100\% au vert} \\
\bottomrule
\end{tabular}
\end{center}

% =====================================================================
\section{Bilan : objectif vs réalisation}
\label{sec:bilan}
% =====================================================================
\subsection{Comparaison au périmètre MVP}
\begin{center}
\begin{tabular}{@{}p{7.6cm}cc@{}}
\toprule
\textbf{Objectif} & \textbf{Prévu (MVP)} & \textbf{Réalisé} \\
\midrule
Chargement de données financières                 & Oui (CSV) & Oui (\textbf{temps réel}) \\
Stratégies d'analyse technique                    & Oui & Oui (MACD, RSI, MA) \\
Simulation des ordres / signaux                   & Oui & Oui \\
Rendement cumulatif                               & Oui & Oui \\
Interface web + graphique + performance           & Oui & Oui (multi-pages) \\
API publique temps réel (yfinance)                & Après-MVP & \textbf{Réalisé} \\
Base de données / historique utilisateur          & Après-MVP & \textbf{Réalisé} \\
Métriques de risque avancées (Sharpe, drawdown)   & Après-MVP & \textbf{Réalisé} \\
Authentification (Google OAuth)                    & Non prévu & \textbf{Réalisé} \\
\bottomrule
\end{tabular}
\end{center}
Le projet a donc non seulement atteint les objectifs du MVP, mais a aussi livré
l'essentiel des fonctionnalités « après-MVP » envisagées, ce qui constitue un
résultat très satisfaisant au regard du périmètre initial.

\subsection{Limitations connues et dette technique}
Par souci de transparence, nous recensons les limites identifiées, pistes
d'amélioration pour la suite :
\begin{itemize}[leftmargin=1.4em]
  \item \textbf{Sécurité / autorisation} : les routes d'historique et de gestion de
        compte ne vérifient pas encore la propriété des ressources côté serveur
        (l'identité est gérée principalement côté client). Un contrôle
        d'autorisation serveur (vérification de session/jeton) est à renforcer.
  \item \textbf{Paramètres financiers} : le capital initial et le \emph{stop-loss}
        sont saisis et stockés mais pas encore intégrés au calcul ; les frais de
        transaction et le \emph{slippage} ne sont pas modélisés.
  \item \textbf{Indicateurs} : le RSI utilise une moyenne exponentielle (approche
        simplifiée vs. RSI de Wilder) ; MACD et MA sont toujours « en position ».
  \item \textbf{Robustesse} : peu de gestion d'erreurs autour des appels réseau
        (yfinance) et de l'écriture en base.
\end{itemize}

\subsection{Perspectives}
Renforcement de l'autorisation serveur, prise en compte du stop-loss et des frais
dans le moteur, ajout de stratégies et de métriques, sauvegarde de stratégies et
listes de suivi (tables déjà modélisées), et amélioration de la robustesse réseau.

% =====================================================================
\section{Manuel utilisateur}
% =====================================================================
\subsection{Prérequis}
Python~3.11+, Node.js~20+ et npm~10+.

\subsection{Installation et lancement}
\paragraph{Backend.}
\begin{lstlisting}[style=sh]
cd backend
python -m venv .venv
.\.venv\Scripts\Activate.ps1     # Windows (PowerShell)
# source .venv/bin/activate      # Linux / macOS
pip install -r requirements.txt
uvicorn main:app --reload --host 127.0.0.1 --port 8000
\end{lstlisting}
API disponible sur \url{http://127.0.0.1:8000}, documentation interactive sur
\url{http://127.0.0.1:8000/docs}.

\paragraph{Frontend.}
\begin{lstlisting}[style=sh]
cd frontend
npm install
npm start
\end{lstlisting}
Application disponible sur \url{http://localhost:4200}.

\subsection{Utilisation pas à pas}
\figph{screenshot-login.png}{0.6\linewidth}{5.5cm}{Page de connexion : authentification Google ou accès en mode démo.}
\begin{enumerate}[leftmargin=1.6em]
  \item \textbf{Connexion} : sur la page de connexion, se connecter avec Google ou
        cliquer sur « Mode démo » pour utiliser l'application sans compte.
  \item \textbf{Données de marché} (page \emph{Data}) : choisir un indice et une
        période pour afficher le graphique des cours.
  \item \textbf{Backtest} (page \emph{Backtests}) : sélectionner indice, période,
        stratégie et ajuster les paramètres, puis cliquer sur « Lancer ».
  \item \textbf{Résultats} (page \emph{Results}) : consulter les métriques
        (rendement, Sharpe, drawdown, trades, taux de réussite), la courbe d'équité
        et le tableau des signaux ; naviguer dans l'historique récent.
  \item \textbf{Tableau de bord} (page \emph{Dashboard}) : vue synthétique des
        backtests et du meilleur rendement.
  \item \textbf{Paramètres} (page \emph{Settings}) : modifier le nom d'utilisateur,
        ajuster les préférences d'interface, supprimer le compte.
\end{enumerate}

\subsection{Exécution des tests (développeurs)}
\begin{lstlisting}[style=sh]
# Backend
cd backend
pip install -r requirements-dev.txt
pytest                          # 58 tests

# Frontend
cd frontend
npm install
npm test                        # ou : ng test --no-watch  (41 tests)
\end{lstlisting}

% =====================================================================
\section{Conclusion}
% =====================================================================
Ce projet nous a permis de concevoir et de livrer une plateforme web de
backtesting fonctionnelle, dépassant les objectifs du MVP grâce à l'intégration de
données temps réel, d'une base de données, de métriques avancées et d'une
authentification. La démarche Agile, le découpage en pôles et l'outillage
collaboratif ont structuré un développement incrémental. La campagne de tests
(99~tests automatisés, tous au vert) garantit la fiabilité du c\oe{}ur
métier et facilite la maintenance future. Les limitations identifiées constituent
une feuille de route claire pour la poursuite du projet.

% =====================================================================
\section{Références}
% =====================================================================
\begin{itemize}[leftmargin=1.4em]
  \item Angular --- \url{https://angular.dev}
  \item Angular Material --- \url{https://material.angular.io}
  \item lightweight-charts (TradingView) --- \url{https://tradingview.github.io/lightweight-charts/}
  \item FastAPI --- \url{https://fastapi.tiangolo.com/}
  \item Pydantic --- \url{https://docs.pydantic.dev/}
  \item Pandas --- \url{https://pandas.pydata.org/docs/}
  \item yfinance --- \url{https://github.com/ranaroussi/yfinance}
  \item SQLAlchemy --- \url{https://docs.sqlalchemy.org/}
  \item Authlib --- \url{https://docs.authlib.org/}
  \item pytest --- \url{https://docs.pytest.org/}
  \item Vitest --- \url{https://vitest.dev/}
  \item Investopedia --- Backtesting --- \url{https://www.investopedia.com/terms/b/backtesting.asp}
  \item Formation Git : D. Ranc, « Introduction à GIT pour TSP1A », PRO3600.
  \item Méthode Agile : H. Charbonnier, PRO3100 « GATE ».
\end{itemize}

% =====================================================================
\section{Lexique}
% =====================================================================
\begin{description}[leftmargin=1.2em,style=nextline]
  \item[API] Interface de programmation permettant au frontend de communiquer avec
        le backend (échanges en JSON).
  \item[Backtesting] Simulation d'une stratégie d'investissement sur des données
        historiques pour en évaluer la viabilité.
  \item[OHLCV] Format standard de données financières : Open, High, Low, Close,
        Volume.
  \item[Drawdown (max)] Pire perte depuis un sommet de la courbe d'équité.
  \item[Ratio de Sharpe] Rendement ajusté du risque (rendement moyen rapporté à la
        volatilité), ici annualisé.
  \item[MACD] \emph{Moving Average Convergence Divergence} : indicateur fondé sur la
        différence de deux moyennes exponentielles.
  \item[RSI] \emph{Relative Strength Index} : oscillateur de momentum borné entre 0
        et 100.
  \item[MA Crossover] Stratégie de croisement de moyennes mobiles (rapide/lente).
  \item[Signal] Position décidée par la stratégie : $+1$ (long), $-1$ (court),
        $0$ (neutre).
  \item[Courbe d'équité] Évolution de la valeur cumulée d'un portefeuille au cours
        du temps.
  \item[Look-ahead bias] Biais consistant à utiliser une information future ;
        évité ici par le décalage du signal.
  \item[OAuth] Protocole d'autorisation déléguée (ici via Google).
  \item[ORM] \emph{Object-Relational Mapping} : correspondance objets/tables
        (SQLAlchemy).
  \item[SPA] \emph{Single Page Application} : application web ne rechargeant qu'une
        seule page.
  \item[MVP] \emph{Produit Minimum Viable} : version minimale validant le concept.
  \item[Sprint] Cycle de développement court et itératif (méthode Agile).
\end{description}

\end{document}
